% Options for packages loaded elsewhere
\PassOptionsToPackage{unicode}{hyperref}
\PassOptionsToPackage{hyphens}{url}
\documentclass[
]{article}
\usepackage{xcolor}
\usepackage[margin=1in]{geometry}
\usepackage{amsmath,amssymb}
\setcounter{secnumdepth}{-\maxdimen} % remove section numbering
\usepackage{iftex}
\ifPDFTeX
  \usepackage[T1]{fontenc}
  \usepackage[utf8]{inputenc}
  \usepackage{textcomp} % provide euro and other symbols
\else % if luatex or xetex
  \usepackage{unicode-math} % this also loads fontspec
  \defaultfontfeatures{Scale=MatchLowercase}
  \defaultfontfeatures[\rmfamily]{Ligatures=TeX,Scale=1}
\fi
\usepackage{lmodern}
\ifPDFTeX\else
  % xetex/luatex font selection
\fi
% Use upquote if available, for straight quotes in verbatim environments
\IfFileExists{upquote.sty}{\usepackage{upquote}}{}
\IfFileExists{microtype.sty}{% use microtype if available
  \usepackage[]{microtype}
  \UseMicrotypeSet[protrusion]{basicmath} % disable protrusion for tt fonts
}{}
\makeatletter
\@ifundefined{KOMAClassName}{% if non-KOMA class
  \IfFileExists{parskip.sty}{%
    \usepackage{parskip}
  }{% else
    \setlength{\parindent}{0pt}
    \setlength{\parskip}{6pt plus 2pt minus 1pt}}
}{% if KOMA class
  \KOMAoptions{parskip=half}}
\makeatother
\usepackage{color}
\usepackage{fancyvrb}
\newcommand{\VerbBar}{|}
\newcommand{\VERB}{\Verb[commandchars=\\\{\}]}
\DefineVerbatimEnvironment{Highlighting}{Verbatim}{commandchars=\\\{\}}
% Add ',fontsize=\small' for more characters per line
\usepackage{framed}
\definecolor{shadecolor}{RGB}{248,248,248}
\newenvironment{Shaded}{\begin{snugshade}}{\end{snugshade}}
\newcommand{\AlertTok}[1]{\textcolor[rgb]{0.94,0.16,0.16}{#1}}
\newcommand{\AnnotationTok}[1]{\textcolor[rgb]{0.56,0.35,0.01}{\textbf{\textit{#1}}}}
\newcommand{\AttributeTok}[1]{\textcolor[rgb]{0.13,0.29,0.53}{#1}}
\newcommand{\BaseNTok}[1]{\textcolor[rgb]{0.00,0.00,0.81}{#1}}
\newcommand{\BuiltInTok}[1]{#1}
\newcommand{\CharTok}[1]{\textcolor[rgb]{0.31,0.60,0.02}{#1}}
\newcommand{\CommentTok}[1]{\textcolor[rgb]{0.56,0.35,0.01}{\textit{#1}}}
\newcommand{\CommentVarTok}[1]{\textcolor[rgb]{0.56,0.35,0.01}{\textbf{\textit{#1}}}}
\newcommand{\ConstantTok}[1]{\textcolor[rgb]{0.56,0.35,0.01}{#1}}
\newcommand{\ControlFlowTok}[1]{\textcolor[rgb]{0.13,0.29,0.53}{\textbf{#1}}}
\newcommand{\DataTypeTok}[1]{\textcolor[rgb]{0.13,0.29,0.53}{#1}}
\newcommand{\DecValTok}[1]{\textcolor[rgb]{0.00,0.00,0.81}{#1}}
\newcommand{\DocumentationTok}[1]{\textcolor[rgb]{0.56,0.35,0.01}{\textbf{\textit{#1}}}}
\newcommand{\ErrorTok}[1]{\textcolor[rgb]{0.64,0.00,0.00}{\textbf{#1}}}
\newcommand{\ExtensionTok}[1]{#1}
\newcommand{\FloatTok}[1]{\textcolor[rgb]{0.00,0.00,0.81}{#1}}
\newcommand{\FunctionTok}[1]{\textcolor[rgb]{0.13,0.29,0.53}{\textbf{#1}}}
\newcommand{\ImportTok}[1]{#1}
\newcommand{\InformationTok}[1]{\textcolor[rgb]{0.56,0.35,0.01}{\textbf{\textit{#1}}}}
\newcommand{\KeywordTok}[1]{\textcolor[rgb]{0.13,0.29,0.53}{\textbf{#1}}}
\newcommand{\NormalTok}[1]{#1}
\newcommand{\OperatorTok}[1]{\textcolor[rgb]{0.81,0.36,0.00}{\textbf{#1}}}
\newcommand{\OtherTok}[1]{\textcolor[rgb]{0.56,0.35,0.01}{#1}}
\newcommand{\PreprocessorTok}[1]{\textcolor[rgb]{0.56,0.35,0.01}{\textit{#1}}}
\newcommand{\RegionMarkerTok}[1]{#1}
\newcommand{\SpecialCharTok}[1]{\textcolor[rgb]{0.81,0.36,0.00}{\textbf{#1}}}
\newcommand{\SpecialStringTok}[1]{\textcolor[rgb]{0.31,0.60,0.02}{#1}}
\newcommand{\StringTok}[1]{\textcolor[rgb]{0.31,0.60,0.02}{#1}}
\newcommand{\VariableTok}[1]{\textcolor[rgb]{0.00,0.00,0.00}{#1}}
\newcommand{\VerbatimStringTok}[1]{\textcolor[rgb]{0.31,0.60,0.02}{#1}}
\newcommand{\WarningTok}[1]{\textcolor[rgb]{0.56,0.35,0.01}{\textbf{\textit{#1}}}}
\usepackage{graphicx}
\makeatletter
\newsavebox\pandoc@box
\newcommand*\pandocbounded[1]{% scales image to fit in text height/width
  \sbox\pandoc@box{#1}%
  \Gscale@div\@tempa{\textheight}{\dimexpr\ht\pandoc@box+\dp\pandoc@box\relax}%
  \Gscale@div\@tempb{\linewidth}{\wd\pandoc@box}%
  \ifdim\@tempb\p@<\@tempa\p@\let\@tempa\@tempb\fi% select the smaller of both
  \ifdim\@tempa\p@<\p@\scalebox{\@tempa}{\usebox\pandoc@box}%
  \else\usebox{\pandoc@box}%
  \fi%
}
% Set default figure placement to htbp
\def\fps@figure{htbp}
\makeatother
\setlength{\emergencystretch}{3em} % prevent overfull lines
\providecommand{\tightlist}{%
  \setlength{\itemsep}{0pt}\setlength{\parskip}{0pt}}
\usepackage{bookmark}
\IfFileExists{xurl.sty}{\usepackage{xurl}}{} % add URL line breaks if available
\urlstyle{same}
\hypersetup{
  pdftitle={Calibration{,} discrimination{,} ROC/AUC{,} Brier score},
  hidelinks,
  pdfcreator={LaTeX via pandoc}}

\title{Calibration, discrimination, ROC/AUC, Brier score}
\author{}
\date{\vspace{-2.5em}}

\begin{document}
\maketitle

\textbf{Inference labs} use the five-step template: Hypothesis →
Visualise → Assumptions → Conduct → Conclude.

\subsection{Learning objectives}\label{learning-objectives}

\begin{itemize}
\tightlist
\item
  Separate discrimination (AUC) from calibration and from overall
  performance (Brier score).
\item
  Draw an ROC curve with \texttt{pROC} and a calibration plot by binning
  predicted probabilities.
\item
  Understand that a high AUC does not imply well-calibrated
  probabilities.
\end{itemize}

\subsection{Prerequisites}\label{prerequisites}

Logistic regression.

\subsection{Background}\label{background}

A risk prediction model should do three things well: separate events
from non-events (discrimination), produce predicted probabilities that
match observed frequencies (calibration), and balance both in a single
summary of overall performance (the Brier score). A model can have high
discrimination and poor calibration --- for instance, if the predicted
probabilities are systematically too extreme --- and a model with
perfect calibration can still be useless if it cannot separate cases
from non-cases.

ROC analysis walks the decision threshold across all possible values and
plots sensitivity against 1 − specificity. The area under the curve
(AUC) is the probability that a randomly chosen case receives a higher
predicted probability than a randomly chosen non-case. The Brier score
is the mean squared error between predicted probabilities and the 0/1
outcome; a perfect model has a Brier of 0 and the no-skill model has
Brier equal to the variance of the outcome.

Calibration plots group predicted probabilities into bins and plot the
observed event rate in each bin against the mean predicted probability.
A 45° reference line is perfect calibration. Smoothed (LOESS)
calibration curves are an alternative in larger samples.

\subsection{Setup}\label{setup}

\begin{Shaded}
\begin{Highlighting}[]
\FunctionTok{library}\NormalTok{(tidyverse)}
\FunctionTok{library}\NormalTok{(broom)}
\FunctionTok{library}\NormalTok{(MASS)}
\FunctionTok{library}\NormalTok{(pROC)}
\FunctionTok{set.seed}\NormalTok{(}\DecValTok{42}\NormalTok{)}
\FunctionTok{theme\_set}\NormalTok{(}\FunctionTok{theme\_minimal}\NormalTok{(}\AttributeTok{base\_size =} \DecValTok{12}\NormalTok{))}
\end{Highlighting}
\end{Shaded}

\subsection{1. Hypothesis}\label{hypothesis}

Using \texttt{MASS::Pima.tr} (train) and \texttt{MASS::Pima.te} (test),
fit a logistic regression and evaluate its discrimination, calibration,
and Brier score.

\subsection{2. Visualise}\label{visualise}

\begin{Shaded}
\begin{Highlighting}[]
\NormalTok{test  }\OtherTok{\textless{}{-}} \FunctionTok{as\_tibble}\NormalTok{(Pima.te) }\SpecialCharTok{|\textgreater{}} \FunctionTok{mutate}\NormalTok{(}\AttributeTok{y =} \FunctionTok{as.integer}\NormalTok{(type }\SpecialCharTok{==} \StringTok{"Yes"}\NormalTok{))}

\NormalTok{fit }\OtherTok{\textless{}{-}} \FunctionTok{glm}\NormalTok{(y }\SpecialCharTok{\textasciitilde{}}\NormalTok{ glu }\SpecialCharTok{+}\NormalTok{ bmi }\SpecialCharTok{+}\NormalTok{ age }\SpecialCharTok{+}\NormalTok{ ped, }\AttributeTok{data =}\NormalTok{ train, }\AttributeTok{family =}\NormalTok{ binomial)}
\NormalTok{test}\SpecialCharTok{$}\NormalTok{p }\OtherTok{\textless{}{-}} \FunctionTok{predict}\NormalTok{(fit, test, }\AttributeTok{type =} \StringTok{"response"}\NormalTok{)}

\FunctionTok{ggplot}\NormalTok{(test, }\FunctionTok{aes}\NormalTok{(p, }\AttributeTok{fill =} \FunctionTok{factor}\NormalTok{(y))) }\SpecialCharTok{+}
  \FunctionTok{geom\_histogram}\NormalTok{(}\AttributeTok{position =} \StringTok{"identity"}\NormalTok{, }\AttributeTok{alpha =} \FloatTok{0.5}\NormalTok{, }\AttributeTok{bins =} \DecValTok{30}\NormalTok{) }\SpecialCharTok{+}
  \FunctionTok{labs}\NormalTok{(}\AttributeTok{x =} \StringTok{"Predicted probability"}\NormalTok{, }\AttributeTok{y =} \StringTok{"Count"}\NormalTok{, }\AttributeTok{fill =} \StringTok{"Diabetes"}\NormalTok{)}
\end{Highlighting}
\end{Shaded}

\subsection{3. Assumptions}\label{assumptions}

The usual logistic-regression assumptions for the training fit;
evaluation is on held-out data.

\subsection{4. Conduct}\label{conduct}

Discrimination:

\begin{Shaded}
\begin{Highlighting}[]
\FunctionTok{auc}\NormalTok{(roc\_obj)}
\FunctionTok{ci.auc}\NormalTok{(roc\_obj)}

\FunctionTok{plot}\NormalTok{(roc\_obj, }\AttributeTok{main =} \FunctionTok{sprintf}\NormalTok{(}\StringTok{"ROC — AUC = \%.2f"}\NormalTok{, }\FunctionTok{auc}\NormalTok{(roc\_obj)))}
\end{Highlighting}
\end{Shaded}

Calibration (decile binning):

\begin{Shaded}
\begin{Highlighting}[]
  \FunctionTok{mutate}\NormalTok{(}\AttributeTok{bin =} \FunctionTok{ntile}\NormalTok{(p, }\DecValTok{10}\NormalTok{)) }\SpecialCharTok{|\textgreater{}}
  \FunctionTok{group\_by}\NormalTok{(bin) }\SpecialCharTok{|\textgreater{}}
  \FunctionTok{summarise}\NormalTok{(}\AttributeTok{predicted =} \FunctionTok{mean}\NormalTok{(p), }\AttributeTok{observed =} \FunctionTok{mean}\NormalTok{(y), }\AttributeTok{n =} \FunctionTok{n}\NormalTok{())}

\FunctionTok{ggplot}\NormalTok{(cal, }\FunctionTok{aes}\NormalTok{(predicted, observed)) }\SpecialCharTok{+}
  \FunctionTok{geom\_abline}\NormalTok{(}\AttributeTok{linetype =} \DecValTok{2}\NormalTok{, }\AttributeTok{colour =} \StringTok{"grey50"}\NormalTok{) }\SpecialCharTok{+}
  \FunctionTok{geom\_point}\NormalTok{(}\AttributeTok{size =} \DecValTok{3}\NormalTok{) }\SpecialCharTok{+}
  \FunctionTok{geom\_line}\NormalTok{() }\SpecialCharTok{+}
  \FunctionTok{coord\_equal}\NormalTok{() }\SpecialCharTok{+}
  \FunctionTok{labs}\NormalTok{(}\AttributeTok{x =} \StringTok{"Predicted probability (bin mean)"}\NormalTok{,}
       \AttributeTok{y =} \StringTok{"Observed event rate"}\NormalTok{)}
\end{Highlighting}
\end{Shaded}

Brier score:

\begin{Shaded}
\begin{Highlighting}[]
\NormalTok{brier}
\end{Highlighting}
\end{Shaded}

\subsection{5. Concluding statement}\label{concluding-statement}

\begin{quote}
On the held-out test set (\emph{n} = \texttt{nrow(test)}), the
four-predictor logistic model achieved AUC =
\texttt{round(as.numeric(auc(roc\_obj)),\ 2)} (95\% CI:
\texttt{round(ci.auc(roc\_obj){[}1{]},\ 2)} to
\texttt{round(ci.auc(roc\_obj){[}3{]},\ 2)}) and a Brier score of
\texttt{round(brier,\ 3)}. Calibration by decile was close to the 45°
line.
\end{quote}

\subsection{Common pitfalls}\label{common-pitfalls}

\begin{itemize}
\tightlist
\item
  Reporting AUC on the training set without cross-validation.
\item
  Using accuracy at a 0.5 threshold when the prevalence is not 0.5.
\item
  Treating a single deciles-calibration plot as a formal test.
\end{itemize}

\subsection{Further reading}\label{further-reading}

\begin{itemize}
\tightlist
\item
  Steyerberg EW. \emph{Clinical Prediction Models}.
\item
  Harrell FE. \emph{Regression Modeling Strategies}, ch.~10.
\item
  Van Calster B et al.~(2019), \emph{Calibration: the Achilles
  heel\ldots{}}
\end{itemize}

\subsection{Session info}\label{session-info}

\begin{Shaded}
\begin{Highlighting}[]

\end{Highlighting}
\end{Shaded}

\subsection{See also --- chapter index}\label{see-also-chapter-index}

\begin{itemize}
\tightlist
\item
  \href{../index.Rmd}{Methods in this chapter}
\item
  \href{../../../ch00-find-your-method.Rmd}{Find your method (Chapter
  0)}
\end{itemize}

\end{document}
